%% OTP_Project_Narrative.tex
%% FY2026 Open Textbook Pilot (84.116T) -- Project Narrative
%% Recommended limit: 65 pages (TOC, budget, abstract, CVs, bibliography,
%% logic model, and letters of support do NOT count toward the limit).
%% The TOC must not exceed ONE double-spaced page.
\documentclass[12pt,letterpaper]{article}
\usepackage{../otpgrant}

\title{\bfseries Project Narrative\\[0.5ex]
  \large Ximera: Interactive Mathematics Education Resources for All\\
  Open Textbook Pilot Program (84.116T), FY 2026}
\author{The Ohio State University}
\date{}

\begin{document}
\doublespacing % ED standard for the narrative

\maketitle
\thispagestyle{empty}
\clearpage

\pagenumbering{roman}
\tableofcontents % must fit on one double-spaced page
\clearpage
\pagenumbering{arabic}

%% =====================================================================
%% The notice requires every priority being addressed to be identified
%% in BOTH the abstract and the project narrative. The structure below
%% follows the notice in order: AP1, AP2, AP3, then selection criteria
%% (a)-(e), so reviewers can score as they read.
%% =====================================================================

\section{Introduction}

\subsection{What is Ximera?}


\subsection{Current state of the project}

Ximera, pronounced ``chimera,'' (\textbf{X}imera:
\textbf{I}nteractive, \textbf{M}athematics, \textbf{E}ducation,
\textbf{R}esources, for \textbf{A}ll) is an open-source platform that
provides tools for authoring and publishing (PDF and Online),
open-source, interactive educational content, such as
textbooks, assessments, and online courses.  
\textbf{The ultimate goal of this project is to promote sustained student success and savings.}


Students interact with the free \textit{content} produced within
Ximera, hence their experience is highly dependent on the \textit{quality} of this
    content. Research \ref{F21} shows that \textbf{students find Ximera materials to be more readable
    than traditional course materials and perform equivalently to those using
    proprietary textbooks and online homework systems}. While students typically
    encounter Ximera through their courses, many discover it via web-search and
    use the platform as independent learners. In 2025, Ximera servers had over one million
    unique visitors. Since Ximera materials are free, they are accessible to
    anyone, regardless of enrollment in official courses.

Instructors (who are not authors) can freely use
    any Ximera materials,
    \textbf{without permission}, by simply sharing the URLs of activities
    with their students. Many first time Ximera instructors start  by
    using online Ximera content as
    ungraded assignments. Visit our website for an incomplete list of
    Ximera courses that have been deployed online.

    


Authors write and store their content on their own
    machines and GitHub repositories.
    \textbf{Authors own their content and decide how to license their content}. From a
    single source written in \LaTeX, Ximera generates various output: PDF
    worksheets,
    PDF textbooks, and	PDF solution manuals, and so on. Of most interest,
    Ximera can
    also create online interactive activities:
    \begin{center}
        \begin{tikzpicture}
            \node at (-1.8,.2)
            {\resizebox{.65cm}{!}{\input{document.tex}}};
            \node at (-1.8,.2) {\small PDF};
            \node at (-2,0) {\resizebox{.65cm}{!}{\input{document.tex}}};
            \node at (-2,0) {\small PDF};
            \node at (-2.2,-.2)
            {\resizebox{.65cm}{!}{\input{document.tex}}};
            \node at (-2.2,-.2) {\small PDF};
            \draw[->] (-.6,0) -- (-1.4,0);
            \draw[->] (-.6,.2) -- (-1.4,.2);
            \draw[->] (-.6,-.2) -- (-1.4,-.2);
            \node at (0,0) {\resizebox{1cm}{!}{\input{document.tex}}};
            \node at (0,0) {\LaTeX};
            \draw[->] (.6,0) -- (1.4,0);

            \node at (2,0) {\resizebox{1cm}{!}{\input{server.tex}}};
            \draw[->] (2.6,0) -- (3.4,0);
            \node at (4,0) {\resizebox{1cm}{!}{\input{computer.tex}}};
            \node at (-2,-1) {Various};
            \node at (-2,-1.4) {PDFs};
            \node at (0,-1) {Single};
            \node at (0,-1.4) {Source};
            \node at (2,-1) {Deploy};
            \node at (2,-1.4) {Server};
            \node at (4,-1) {Students};
            \node at (4,-1.4) {Engage};

        \end{tikzpicture}
    \end{center}
    The source code used to produce PDFs can also create interactive online activities when deployed to a Ximera server. Students access this content via a URL or an assignment in their LMS.


\section{Absolute Priority 1: Improving Collaboration and Dissemination}
% Consortium of >= 3 IHEs beyond informal partners, plus employers,
% workforce stakeholders, and/or trade or professional associations.
% Address: (a) reducing cost of college for large numbers of students;
% (b) aligning learning objectives with skills/knowledge needed by
% large numbers of students (or industry standards for CTE programs).

\section{Absolute Priority 2: Addressing Gaps in the Open Textbook Marketplace and Bringing Solutions to Scale}
% Comprehensive plan covering all six elements:
\subsection{Assessment of Existing Open Educational Resources}   % (a)
\subsection{Creation and Expansion of Scalable Materials}        % (b)
\subsection{Review Protocols: Accuracy, Rigor, and Accessibility}% (c)
\subsection{Dissemination and Adoption Beyond the Consortium}    % (d)
\subsection{Professional Development for Faculty and Staff}      % (e)
\subsection{Courses for Which Materials Will Be Developed}       % (f)

\section{Absolute Priority 3: Promoting Student Success}
% High-enrollment courses/programs; highest level of student savings.
\subsection{Promoting and Tracking Use; Projected Direct Cost Savings} % (a)
\subsection{Monitoring Impact on Instruction, Outcomes, and Costs}     % (b)
\subsection{Evidence-Based Practices and Dissemination}                % (c)
\subsection{Sustaining and Updating Materials Beyond the Funded Period}% (d)

%% =====================================================================
%% Selection criteria (34 CFR 75.210) -- 100 points total
%% =====================================================================

\criterion{Need for the Project}{12}
\subsection{Data Demonstrating the Issue, Challenge, or Opportunity} % up to 6
\subsection{Nature and Magnitude of Gaps and How the Project Addresses Them} % up to 6

\criterion{Quality of the Project Design}{48}
\subsection{Community Engagement and Continuous Improvement}            % up to 4
\subsection{Clear, Measurable, Ambitious yet Achievable Goals}          % up to 4
\subsection{Specific, Measurable Targets and Reliable Administrative Data} % up to 4
\subsection{Building on Ideas and Efforts from Similar External Projects}  % up to 4
\subsection{Quality, Intensity, and Duration of Professional Development}  % up to 4
\subsection{Meaningful Improvements in Student Achievement}             % up to 4
\subsection{Skills and Competencies for High-Quality Employment}        % up to 4
\subsection{Efficient Strategies, Technology, and Leveraged Resources}  % up to 4
\subsection{Capacity to Bring the Project to Scale}                     % up to 10
\subsection{Mechanisms for Broad Dissemination}                         % up to 6

\criterion{Adequacy of Resources}{12}
\subsection{Relevance and Commitment of Each Partner}     % up to 6
\subsection{Institutionalization Beyond Federal Funding}  % up to 6

\criterion{Quality of the Management Plan}{16}
\subsection{Feasibility: Responsibilities, Timelines, and Milestones} % up to 8
\subsection{Ensuring High-Quality, Accessible Products and Services}  % up to 8

\criterion{Quality of the Evaluation Plan}{12}
\subsection{Thorough, Feasible, Relevant, and Appropriate Methods}    % up to 6
\subsection{Analytic Strategy Linking Components, Mediators, and Outcomes} % up to 6

\end{document}
